\documentclass[11pt,a4paper]{article}
\usepackage[top=2.3cm,bottom=2.3cm,left=2.2cm,right=2.2cm]{geometry}

\usepackage{amsmath,amssymb}
\usepackage{fontspec}
\setmainfont{Liberation Serif}
\setsansfont{Liberation Sans}
\setmonofont{DejaVu Sans Mono}[Scale=0.84]

\usepackage{xcolor}
\definecolor{brandblue}{RGB}{20, 55, 105}
\definecolor{accentblue}{RGB}{35, 95, 165}
\definecolor{codebg}{RGB}{248, 249, 251}
\definecolor{codeframe}{RGB}{215, 222, 230}
\definecolor{javakeyword}{RGB}{0, 45, 170}
\definecolor{javastring}{RGB}{5, 120, 20}
\definecolor{javacomment}{RGB}{120, 120, 120}
\definecolor{javanumber}{RGB}{20, 75, 220}

\definecolor{noteborder}{RGB}{30, 110, 65}
\definecolor{notebg}{RGB}{242, 249, 245}
\definecolor{warnborder}{RGB}{185, 75, 10}
\definecolor{warnbg}{RGB}{254, 248, 240}
\definecolor{tipborder}{RGB}{30, 95, 160}
\definecolor{tipbg}{RGB}{242, 247, 254}

\usepackage{fancyhdr}
\usepackage{lastpage}
\pagestyle{fancy}
\fancyhf{}
\lhead{\parbox[t]{0.58\textwidth}{\raggedright\small\sffamily\color{brandblue}\textbf{T14 --- Classi Annidate, Interne, Locali ed Anonime}}}
\rhead{\small\sffamily\color{gray}Programmazione II -- UniVR (2025/2026)}
\cfoot{\small\sffamily Pagina \thepage\ di \pageref{LastPage}}
\renewcommand{\headrulewidth}{0.5pt}
\renewcommand{\headrule}{\hbox to\headwidth{\color{brandblue!70!black}\leaders\hrule height \headrulewidth\hfill}}

\usepackage{titlesec}
\titleformat{\section}{\Large\bfseries\sffamily\color{brandblue}}{\thesection}{0.8em}{}[\titlerule]
\titleformat{\subsection}{\large\bfseries\sffamily\color{accentblue}}{\thesubsection}{0.8em}{}
\titleformat{\subsubsection}{\normalsize\bfseries\sffamily\color{brandblue!85!black}}{\thesubsubsection}{0.8em}{}

\usepackage{needspace}
\usepackage{fancyvrb}
\DefineVerbatimEnvironment{asciibox}{Verbatim}{
  fontsize=\fontsize{7.5pt}{9.2pt}\selectfont,
  frame=single,
  rulecolor=\color{codeframe},
  fillcolor=\color{codebg},
  framesep=2.5mm
}
\usepackage{listings}
\lstset{
  backgroundcolor=\color{codebg},
  frame=single,
  rulecolor=\color{codeframe},
  basicstyle=\ttfamily\footnotesize,
  keywordstyle=\color{javakeyword}\bfseries,
  stringstyle=\color{javastring},
  commentstyle=\color{javacomment}\itshape,
  numberstyle=\tiny\color{gray},
  numbers=left,
  numbersep=7pt,
  stepnumber=1,
  breaklines=true,
  breakatwhitespace=true,
  showstringspaces=false,
  tabsize=4,
  xleftmargin=12pt,
  framexleftmargin=12pt
}

\newcommand{\passthrough}[1]{#1}
\providecommand{\tightlist}{\setlength{\itemsep}{0pt}\setlength{\parskip}{0pt}}

\usepackage{tcolorbox}
\tcbuselibrary{skins,breakable}
\newtcolorbox{studynote}[1][Nota]{
  enhanced,
  breakable,
  colback=notebg,
  colframe=noteborder,
  fonttitle=\bfseries\sffamily\small,
  coltitle=white,
  title={#1},
  arc=2mm,
  boxrule=0.8pt,
  left=9pt, right=9pt, top=7pt, bottom=7pt
}

\newtcolorbox{studywarning}[1][Attenzione]{
  enhanced,
  breakable,
  colback=warnbg,
  colframe=warnborder,
  fonttitle=\bfseries\sffamily\small,
  coltitle=white,
  title={#1},
  arc=2mm,
  boxrule=0.8pt,
  left=9pt, right=9pt, top=7pt, bottom=7pt
}

\newtcolorbox{studytip}[1][Suggerimento]{
  enhanced,
  breakable,
  colback=tipbg,
  colframe=tipborder,
  fonttitle=\bfseries\sffamily\small,
  coltitle=white,
  title={#1},
  arc=2mm,
  boxrule=0.8pt,
  left=9pt, right=9pt, top=7pt, bottom=7pt
}

\usepackage{booktabs}
\usepackage{longtable}
\usepackage{array}
\usepackage{calc}

\usepackage{hyperref}
\hypersetup{
  colorlinks=true,
  linkcolor=accentblue,
  urlcolor=accentblue,
  citecolor=brandblue,
  pdfborder={0 0 0}
}

\title{\Huge\bfseries\sffamily\color{brandblue} T14 --- Classi Annidate, Interne, Locali ed Anonime \\[0.3em] \large\sffamily Static Nested vs Inner Classes, Enclosing Pointer, Shadowing, Effectively Final e Classi Anonime}
\author{\textbf{Docente:} Prof. Michele Pasqua \quad---\quad \textbf{Appunti Revisione:} Wally}
\date{A.A. 2025/2026 -- Universit\`a degli Studi di Verona}

\begin{document}
\maketitle
\thispagestyle{fancy}
\vspace{-1.2em}
\noindent\textcolor{brandblue}{\rule{\textwidth}{0.8pt}}
\vspace{1.2em}

\begin{center}\rule{0.5\linewidth}{0.5pt}\end{center}

\subsection{1. Analisi Teorica Approfondita (Slide
T14)}\label{analisi-teorica-approfondita-slide-t14}

La lezione 14 esplora la gerarchia delle classi definite all'interno di
altre classi (meccanismo introdotto fin da Java 1.1), le classi anonime,
l'evoluzione verso le \textbf{espressioni Lambda} (Java 8) e la feature
dei \textbf{Varargs}.

\needspace{7cm}
\begin{asciibox}
                 ┌───────────────────────────┐
                 │  Classi Annidate in Java  │
                 └─────────────┬─────────────┘
                               │
            ┌──────────────────┴──────────────────┐
            ▼                                     ▼
 ┌───────────────────────┐             ┌───────────────────────┐
 │ Static Nested Classes │             │     Inner Classes     │
 │   (static class ...)  │             │     (non-static)      │
 └───────────────────────┘             └───────────┬───────────┘
                                                   │
                  ┌────────────────────────────────┼─────────────────┐
                  ▼                                ▼                 ▼
       ┌──────────────────┐ ┌──────────────────┐ ┌──────────────────┐
       │Member Inner Class│ │   Local Class    │ │ Anonymous Class  │
       │(a livello classe)│ │(dentro un metodo)│ │(senza nome/corpo)│
       └──────────────────┘ └──────────────────┘ └────────┬─────────┘
                                                           │ (SAM Interface)
                                                           ▼
                                                 ┌──────────────────┐
                                                 │Lambda Expression │
                                                 │  (arg -> body)   │
                                                 └──────────────────┘
\end{asciibox}

\begin{center}\rule{0.5\linewidth}{0.5pt}\end{center}

\subsubsection{\texorpdfstring{1.1 Static Nested Classes
(\texttt{static\ class\ Nested})}{1.1 Static Nested Classes (static class Nested)}}\label{static-nested-classes-static-class-nested}

Una classe dichiarata con il modificatore
\passthrough{\lstinline!static!} all'interno di un'altra classe
contenitore (\emph{enclosing class}):

\begin{lstlisting}[language=Java]
package pkg;
class Outer {
    static class Nested {
        // ...
    }
}
\end{lstlisting}

\begin{itemize}
\item
  \textbf{Natura}: È a tutti gli effetti una classe di primo livello
  (\emph{top-level}), ma inserita nel \emph{namespace} della classe che
  la ospita.
\item
  \textbf{Nome completo (Fully Qualified Name)}:
  \passthrough{\lstinline!pkg.Outer.Nested!} (a livello di bytecode su
  disco viene generato il file
  \passthrough{\lstinline!Outer$Nested.class!}).
\item
  \textbf{Accesso ai membri}:

  \begin{itemize}
  \tightlist
  \item
    Trattandosi di un membro statico, \textbf{ha accesso a tutti i
    membri statici} della classe contenitore, \textbf{compresi quelli
    privati (\passthrough{\lstinline!private!})}!
  \item
    \textbf{NON ha accesso} ai membri di istanza (campi o metodi non
    statici) di \passthrough{\lstinline!Outer!}, poiché non esiste alcun
    riferimento implicito ad un'istanza di
    \passthrough{\lstinline!Outer!}.
  \end{itemize}
\item
  \textbf{Istanziazione}: Può essere istanziata in modo totalmente
  indipendente dall'esistenza di un oggetto
  \passthrough{\lstinline!Outer!}:

\begin{lstlisting}[language=Java]
Outer.Nested nestedObject = new Outer.Nested();
\end{lstlisting}
\item
  \textbf{Caso d'uso tipico}: Il pattern \textbf{Builder} (visto in
  \href{file:///home/wally/Documenti/GitHub/Programmazione-II/25-26/Teoria-e-Esercitazioni/Codice-20260902/Lezione14/MavenDate/src/main/java/it/oop/core/Date.java\#L85-L106}{Date.java}),
  classi di utilità, nodi interni di strutture dati (es.
  \passthrough{\lstinline!static class Node<E>!} in un albero o grafo).
\end{itemize}

\begin{center}\rule{0.5\linewidth}{0.5pt}\end{center}

\subsubsection{\texorpdfstring{1.2 Member Inner Classes
(\texttt{class\ Inner}
non-statica)}{1.2 Member Inner Classes (class Inner non-statica)}}\label{member-inner-classes-class-inner-non-statica}

Una classe dichiarata all'interno di un'altra classe \textbf{senza} la
parola chiave \passthrough{\lstinline!static!}:

\begin{lstlisting}[language=Java]
package pkg;
class Outer {
    class Inner {
        // ...
    }
}
\end{lstlisting}

\begin{itemize}
\item
  \textbf{Natura}: Un'istanza di \passthrough{\lstinline!Inner!}
  \textbf{non può esistere senza un'istanza associata di
  \passthrough{\lstinline!Outer!}}.
\item
  \textbf{Riferimento implicito (\passthrough{\lstinline!Outer.this!} o
  \passthrough{\lstinline!this$0!})}:

  \begin{itemize}
  \tightlist
  \item
    Ogni oggetto \passthrough{\lstinline!Inner!} memorizza internamente
    un puntatore sintetico nascosto all'oggetto
    \passthrough{\lstinline!Outer!} che l'ha generato.
  \item
    Tramite questo puntatore, \passthrough{\lstinline!Inner!} ha accesso
    a \textbf{tutti} i membri di \passthrough{\lstinline!Outer!}, sia
    statici che di istanza, inclusi quelli
    \passthrough{\lstinline!private!}!
  \end{itemize}
\item
  \textbf{Sintassi di istanziazione da codice esterno}:

\begin{lstlisting}[language=Java]
Outer outerObject = new Outer();
Outer.Inner innerObject = outerObject.new Inner(); // Sintassi specifica con outerObject.new
\end{lstlisting}
\end{itemize}

\paragraph{\texorpdfstring{Confronto: \texttt{Nested} vs \texttt{Inner}
(Slide
7)}{Confronto: Nested vs Inner (Slide 7)}}\label{confronto-nested-vs-inner-slide-7}

\begin{lstlisting}[language=Java]
class Outer {
    private static int I = 5;
    private int i;

    static class Nested {
        // NON ha accesso al campo di istanza i!
        static void print() { System.out.println(I); } // OK: accede a membro statico privato I
    }

    class Inner {
        // Ha accesso sia a I che al campo di istanza i della sua Outer associata
        void print() { System.out.println(i); }
    }

    public void outerMethod() {
        Nested.print();
        Inner innerObject = new Inner(); // Qui 'this.new Inner()' è implicito
        innerObject.print();
    }
}
\end{lstlisting}

\begin{center}\rule{0.5\linewidth}{0.5pt}\end{center}

\subsubsection{1.3 Local Classes (Classi Locali a un
Metodo)}\label{local-classes-classi-locali-a-un-metodo}

Una classe definita all'interno del corpo di un blocco di codice
(tipicamente un metodo):

\begin{lstlisting}[language=Java]
public void outerMethod() {
    final int j = 1;

    class Local { // NESSUN modificatore di visibilità ammesso (no public, private, protected)
        int addOne() { return j + 1; }
    }

    Local localObject = new Local();
    System.out.println(localObject.addOne()); // Stampa 2
}
\end{lstlisting}

\begin{itemize}
\tightlist
\item
  \textbf{Regole di visibilità}: La classe è visibile e istanziabile
  \textbf{esclusivamente all'interno del blocco/metodo} in cui è
  dichiarata.
\item
  \textbf{Cattura delle variabili locali (\emph{Variable Capture} \&
  \emph{Closures})}:

  \begin{itemize}
  \tightlist
  \item
    Può accedere alle variabili locali e ai parametri del metodo
    contenitore \textbf{solo se sono dichiarati
    \passthrough{\lstinline!final!} oppure sono \emph{effectively
    final}} (ovvero il loro valore non viene mai modificato dopo
    l'inizializzazione).
  \item
    \textbf{Perché questo vincolo severo? (Domanda classica d'esame)}:

    \begin{itemize}
    \tightlist
    \item
      Le variabili locali risiedono nello \textbf{Stack Frame} del
      metodo e vengono distrutte non appena il metodo termina
      (\passthrough{\lstinline!return!}).
    \item
      L'oggetto della classe locale
      (\passthrough{\lstinline!localObject!}) risiede nell'\textbf{Heap}
      e potrebbe sopravvivere ben oltre la fine del metodo (ad esempio
      se restituito come interfaccia o registrato come listener).
    \item
      Per permettere all'oggetto di accedere a
      \passthrough{\lstinline!j!}, il compilatore Java crea un campo
      sintetico nascosto dentro \passthrough{\lstinline!Local!} e vi
      \textbf{copia} il valore di \passthrough{\lstinline!j!} al momento
      della creazione.
    \item
      Se \passthrough{\lstinline!j!} potesse essere riassegnato
      (\passthrough{\lstinline!j++!}), il valore nello stack e la copia
      nell'heap diverrebbero disallineati, rompendo la coerenza del
      linguaggio. Rendendo la variabile immutabile
      (\passthrough{\lstinline!final!}), la copia rimane identica e
      coerente per sempre.
    \end{itemize}
  \end{itemize}
\end{itemize}

\begin{center}\rule{0.5\linewidth}{0.5pt}\end{center}

\subsubsection{\texorpdfstring{1.4 Il Meme della Slide 10: \emph{Objects
inside
objects}}{1.4 Il Meme della Slide 10: Objects inside objects}}\label{il-meme-della-slide-10-objects-inside-objects}

Nella Slide 10, il prof. Pasqua inserisce il celebre meme di Xzibit
(\emph{Pimp My Ride}): \textgreater{} \emph{``YO DAWG, I HEAR YOU LIKE
OBJECT ORIENTED PROGRAMMING SO I NESTED A CLASS INSIDE YOUR CLASS SO YOU
CAN CREATE OBJECTS WHILE YOU CREATE OBJECTS''}

A sottolineare con ironia la tentazione (e il potere) di creare
gerarchie di classi dentro classi per incapsulare comportamenti
strettamente accoppiati.

\begin{center}\rule{0.5\linewidth}{0.5pt}\end{center}

\subsubsection{\texorpdfstring{1.5 Motivazioni Architetturali \& Il Caso
di Studio \emph{Iterator Pattern} (Slide
11-12)}{1.5 Motivazioni Architetturali \& Il Caso di Studio Iterator Pattern (Slide 11-12)}}\label{motivazioni-architetturali-il-caso-di-studio-iterator-pattern-slide-11-12}

Perché annidare le classi? 1. \textbf{Migliore organizzazione del
codice}: collocare classi di supporto logico all'interno del namespace
della sola classe che le utilizza. 2. \textbf{Politica di visibilità
privilegiata (\emph{Privileged Access})}: le classi annidate possono
leggere lo stato privato della classe ospitante senza costringere a
esporre metodi \passthrough{\lstinline!public!} o
\passthrough{\lstinline!package-private!} pericolosi verso l'esterno. 3.
\textbf{Meno complessità nel codebase}: evita la proliferazione di
decine di minuscoli file \passthrough{\lstinline!.java!} separati.

\paragraph{\texorpdfstring{L'Esempio del pattern \texttt{Iterable} /
\texttt{Iterator} (Slide
12)}{L'Esempio del pattern Iterable / Iterator (Slide 12)}}\label{lesempio-del-pattern-iterable-iterator-slide-12}

\begin{itemize}
\item
  \textbf{Senza annidamento (Due classi separate)}:
  \passthrough{\lstinline!CRIterator!} deve essere dichiarata
  all'esterno, riceve l'array \passthrough{\lstinline!Student[] arr!}
  nel costruttore e necessita di visibilità speciale:

\begin{lstlisting}[language=Java]
class ClassRoom implements Iterable {
    private Student[] students;
    @Override
    public Iterator iterator() { return new CRIterator(students); }
}
class CRIterator implements Iterator {
    private Student[] arr;
    private int next = 0;
    CRIterator(Student[] arr) { this.arr = arr; }
    @Override public boolean hasNext() { return next < arr.length; }
    @Override public Object next() { return new Student(arr[next++]); }
}
\end{lstlisting}
\item
  \textbf{Con Inner Class (Soluzione idiomatica Java)}:
  \passthrough{\lstinline!CRIterator!} è una
  \passthrough{\lstinline!Inner Class!} dentro
  \passthrough{\lstinline!ClassRoom!}. Non serve memorizzare un array
  duplicato: accede direttamente al campo privato
  \passthrough{\lstinline!students!} di
  \passthrough{\lstinline!ClassRoom!}!

\begin{lstlisting}[language=Java]
class ClassRoom implements Iterable {
    private Student[] students;
    @Override
    public Iterator iterator() { return new CRIterator(); }

    class CRIterator implements Iterator {
        private int next = 0;
        @Override public boolean hasNext() { return next < students.length; }
        @Override public Object next() { return new Student(students[next++]); }
    }
}
\end{lstlisting}
\end{itemize}

\begin{center}\rule{0.5\linewidth}{0.5pt}\end{center}

\subsubsection{1.6 Anonymous Classes (Classi Anonime, Slide
13-16)}\label{anonymous-classes-classi-anonime-slide-13-16}

Una classe locale \textbf{priva di nome}, definita e istanziata
contestualmente all'interno di una singola espressione. *
\textbf{Vincolo}: Può essere creata solo estendendo una classe esistente
o implementando un'interfaccia: ```java // 1. Estendendo una classe
esistente: Person sam = new Person(``Sam'') \{ @Override public String
toString() \{ return ``It's Sam!''; \} \};

// 2. Implementando un'interfaccia: Time init = new Time() \{ @Override
public int getHours() \{ return 0; \} @Override public int getMinutes()
\{ return 0; \} @Override public int getSeconds() \{ return 0; \} \};
\passthrough{\lstinline"* **Regola aurea sui Costruttori**: **Una classe anonima NON può dichiarare alcun costruttore esplicito!** (Poiché i costruttori devono avere per definizione lo stesso nome della classe, e una classe anonima non ha nome). Eventuali inizializzazioni devono sfruttare gli *instance initializer blocks* `\{ ... \}` o i parametri passati al costruttore della superclasse (`new SuperClass(arg1, arg2) \{ ... \}`). * **Espressioni**: Essendo espressioni, devono terminare con un punto e virgola `;` (se assegnate a una variabile) o possono essere passate direttamente come parametri attuali di un metodo:"}java
tasks.add(new Runnable() \{ @Override public void run() \{
System.out.println(``Running!''); \} \}); ```

\begin{center}\rule{0.5\linewidth}{0.5pt}\end{center}

\subsubsection{1.7 Lambda Expressions (Slide
17-19)}\label{lambda-expressions-slide-17-19}

Introdotte in Java 8, sono una notazione sintattica ultra-compatta per
implementare \textbf{interfacce funzionali} (interfacce con un
\textbf{unico metodo astratto}, note come SAM --- \emph{Single Abstract
Method}).

\begin{lstlisting}
(parametri) -> corpo
\end{lstlisting}

\begin{itemize}
\tightlist
\item
  \textbf{Parametri}:

  \begin{itemize}
  \tightlist
  \item
    Nessun parametro: \passthrough{\lstinline!() -> 42!}
  \item
    Un parametro: \passthrough{\lstinline!x -> x + 1!} (le parentesi
    tonde sono opzionali se non si specifica il tipo esplicito)
  \item
    Due o più parametri: \passthrough{\lstinline!(x, y) -> x + y!}
  \end{itemize}
\item
  \textbf{Corpo}:

  \begin{itemize}
  \tightlist
  \item
    Espressione singola: \passthrough{\lstinline!param -> param - 1!}
    (restituzione implicita del valore calcolato, niente
    \passthrough{\lstinline!return!} né graffe)
  \item
    Blocco di codice:
    \passthrough{\lstinline!param -> \{ int res = param - 1; return res; \}!}
    (richiede graffe esplicite e \passthrough{\lstinline!return!})
  \end{itemize}
\item
  \textbf{Inferenza di tipo (\emph{Type Inference})}: Il compilatore
  Java deduce i tipi dei parametri e del valore di ritorno dal contesto
  d'uso (\emph{target typing} dell'interfaccia funzionale).
\item
  \textbf{Differenza tra Classi Anonime e Lambda}:

  \begin{itemize}
  \tightlist
  \item
    La classe anonima può implementare interfacce con più metodi
    astratti o estendere classi concrete/astratte; la lambda può
    implementare \textbf{solo interfacce funzionali SAM}.
  \item
    La classe anonima crea un nuovo scope per
    \passthrough{\lstinline!this!} (che punta all'istanza anonima
    stessa); nella lambda \passthrough{\lstinline!this!} è
    \textbf{lessicale} (punta all'istanza della classe circostante).
  \item
    La classe anonima può dichiarare campi di istanza che mantengono
    stato; la lambda è per natura \emph{stateless}.
  \end{itemize}
\end{itemize}

\begin{center}\rule{0.5\linewidth}{0.5pt}\end{center}

\subsubsection{\texorpdfstring{1.8 Varargs (\texttt{type...\ args},
Slide
20-21)}{1.8 Varargs (type... args, Slide 20-21)}}\label{varargs-type...-args-slide-20-21}

Consente a un metodo di accettare un numero variabile (da zero a molti)
di argomenti dello stesso tipo:

\begin{lstlisting}[language=Java]
static int min(int... values) {
    int res = Integer.MAX_VALUE;
    for (int v : values)
        if (v < res) res = v;
    return res;
}
\end{lstlisting}

\begin{itemize}
\tightlist
\item
  \textbf{Dietro le quinte del compilatore}: Il parametro
  \passthrough{\lstinline!int... values!} viene tradotto esattamente in
  un array \passthrough{\lstinline!int[] values!}. Al momento della
  chiamata:

  \begin{itemize}
  \tightlist
  \item
    \passthrough{\lstinline!min(1, 0, 5)!} \(\rightarrow\) il
    compilatore inietta il codice:
    \passthrough{\lstinline!min(new int[]\{ 1, 0, 5 \})!}.
  \item
    \passthrough{\lstinline!min()!} (senza argomenti) \(\rightarrow\)
    crea un array vuoto \passthrough{\lstinline!min(new int[]\{\})!}.
  \end{itemize}
\item
  \textbf{Regole tassative}:

  \begin{enumerate}
  \def\labelenumi{\arabic{enumi}.}
  \tightlist
  \item
    Ci può essere \textbf{al massimo un} parametro varargs per metodo.
  \item
    Il parametro varargs deve essere tassativamente l'\textbf{ultimo}
    nella lista dei parametri formali
    (\passthrough{\lstinline!void foo(String s, int... nums)!} è valido;
    \passthrough{\lstinline!void foo(int... nums, String s)!} genera
    errore di compilazione).
  \end{enumerate}
\end{itemize}

\begin{center}\rule{0.5\linewidth}{0.5pt}\end{center}

\subsection{\texorpdfstring{2. Analisi Dettagliata del Codice
(\texttt{Lezione14/MavenDate})}{2. Analisi Dettagliata del Codice (Lezione14/MavenDate)}}\label{analisi-dettagliata-del-codice-lezione14mavendate}

Nella Lezione 14 il progetto
\href{file:///home/wally/Documenti/GitHub/Programmazione-II/25-26/Teoria-e-Esercitazioni/Codice-20260902/Lezione14/MavenDate}{MavenDate}
introduce il primo ponte fondamentale verso i \textbf{Tipi Generici}
(anticipando la teoria formale di Slide T15) e sfrutta le classi
annidate, le classi anonime e le lambda viste nella Lezione 13.

\subsubsection{\texorpdfstring{2.1 Diff di
\href{file:///home/wally/Documenti/GitHub/Programmazione-II/25-26/Teoria-e-Esercitazioni/Codice-20260902/Lezione14/MavenDate/src/main/java/it/oop/core/Date.java\#L1-L80}{Date.java}
(Lezione 13 vs Lezione
14)}{2.1 Diff di Date.java (Lezione 13 vs Lezione 14)}}\label{diff-di-date.java-lezione-13-vs-lezione-14}

Fino alla Lezione 13, \passthrough{\lstinline!Date!} implementava
l'interfaccia non tipizzata (\emph{raw})
\passthrough{\lstinline!Comparable!}:

\begin{lstlisting}
--- Lezione13/Date.java
+++ Lezione14/Date.java
@@ -1,6 +1,6 @@
 package it.oop.core;
 
-public class Date implements Comparable {
+public class Date implements Comparable<Date> {
 
@@ -73,15 +73,14 @@
     @Override
-    public int compareTo(Object other) {
-        Date otherAsDate = (Date) other;
-        int diff = this.year - otherAsDate.getYear();
+    public int compareTo(Date other) {
+        int diff = this.year - other.getYear();
         if (diff != 0)
             return diff;
-        diff = this.month - otherAsDate.getMonth();
+        diff = this.month - other.getMonth();
         if (diff != 0)
             return diff;
-        return this.day - otherAsDate.getDay();
+        return this.day - other.getDay();
     }
\end{lstlisting}

\begin{itemize}
\tightlist
\item
  \textbf{Beneficio immediato del Generics}:

  \begin{enumerate}
  \def\labelenumi{\arabic{enumi}.}
  \tightlist
  \item
    La firma del metodo diventa fortemente tipizzata a compile-time:
    \passthrough{\lstinline!compareTo(Date other)!} anziché
    \passthrough{\lstinline!compareTo(Object other)!}.
  \item
    Viene eliminato il cast esplicito
    \passthrough{\lstinline!(Date) other!}, azzerando il rischio di
    incorrere a runtime in una
    \passthrough{\lstinline!ClassCastException!}.
  \end{enumerate}
\end{itemize}

\begin{center}\rule{0.5\linewidth}{0.5pt}\end{center}

\subsubsection{\texorpdfstring{2.2 La Nuova Classe Generica
\href{file:///home/wally/Documenti/GitHub/Programmazione-II/25-26/Teoria-e-Esercitazioni/Codice-20260902/Lezione14/MavenDate/src/main/java/it/oop/core/Pair.java}{Pair.java}}{2.2 La Nuova Classe Generica Pair.java}}\label{la-nuova-classe-generica-pair.java}

\begin{lstlisting}[language=Java]
package it.oop.core;

public class Pair<F, S> {
    private final F first;
    private final S second;

    public Pair(F first, S second) {
        this.first = first;
        this.second = second;
    }

    public F getFirst() { return first; }
    public S getSecond() { return second; }
}
\end{lstlisting}

\begin{itemize}
\tightlist
\item
  \passthrough{\lstinline!F!} (\emph{First}) e
  \passthrough{\lstinline!S!} (\emph{Second}) sono due \textbf{Type
  Parameters} indipendenti.
\item
  I campi sono \passthrough{\lstinline!private final!}, rendendo la
  coppia immutabile una volta istanziata.
\end{itemize}

\begin{center}\rule{0.5\linewidth}{0.5pt}\end{center}

\subsubsection{\texorpdfstring{2.3 La Classe Generica con Bounded Type
\href{file:///home/wally/Documenti/GitHub/Programmazione-II/25-26/Teoria-e-Esercitazioni/Codice-20260902/Lezione14/MavenDate/src/main/java/it/oop/core/OrderedPair.java}{OrderedPair.java}}{2.3 La Classe Generica con Bounded Type OrderedPair.java}}\label{la-classe-generica-con-bounded-type-orderedpair.java}

\begin{lstlisting}[language=Java]
package it.oop.core;

public class OrderedPair <T extends Comparable<T>> {
    private final T first;
    private final T second;

    public OrderedPair(T first, T second) {
        this.first = first;
        this.second = second;
        if (first.compareTo(second) > 0)
            System.out.println("Pair not ordered!");
    }

    public T getFirst() { return first; }
    public T getSecond() { return second; }
}
\end{lstlisting}

\begin{itemize}
\tightlist
\item
  \textbf{Bounded Type Parameter
  (\passthrough{\lstinline!<T extends Comparable<T>>!})}:

  \begin{itemize}
  \tightlist
  \item
    Il tipo \passthrough{\lstinline!T!} non può essere un qualsiasi tipo
    \passthrough{\lstinline!Object!}: deve tassativamente implementare
    l'interfaccia \passthrough{\lstinline!Comparable<T>!}.
  \item
    Grazie a questo vincolo, il compilatore consente all'interno del
    costruttore di invocare legalmente
    \passthrough{\lstinline!first.compareTo(second)!}.
  \item
    Se si tentasse di istanziare
    \passthrough{\lstinline!OrderedPair<Object>!}, il compilatore
    bloccherebbe l'operazione con un errore prima ancora di eseguire.
  \end{itemize}
\end{itemize}

\begin{center}\rule{0.5\linewidth}{0.5pt}\end{center}

\subsubsection{\texorpdfstring{2.4 Le Sottoclassi Specializzate
\href{file:///home/wally/Documenti/GitHub/Programmazione-II/25-26/Teoria-e-Esercitazioni/Codice-20260902/Lezione14/MavenDate/src/main/java/it/oop/core/DatePair.java}{DatePair.java}
e
\href{file:///home/wally/Documenti/GitHub/Programmazione-II/25-26/Teoria-e-Esercitazioni/Codice-20260902/Lezione14/MavenDate/src/main/java/it/oop/core/DateInterval.java}{DateInterval.java}}{2.4 Le Sottoclassi Specializzate DatePair.java e DateInterval.java}}\label{le-sottoclassi-specializzate-datepair.java-e-dateinterval.java}

\paragraph{\texorpdfstring{\href{file:///home/wally/Documenti/GitHub/Programmazione-II/25-26/Teoria-e-Esercitazioni/Codice-20260902/Lezione14/MavenDate/src/main/java/it/oop/core/DatePair.java}{DatePair.java}}{DatePair.java}}\label{datepair.java}

\begin{lstlisting}[language=Java]
package it.oop.core;

public class DatePair extends Pair<Date, Date> {
    public DatePair(Date left, Date right) {
        super(left, right);
    }

    public Date getLeft() { return getFirst(); }
    public Date getRight() { return getSecond(); }

    @Override
    public String toString() {
        Date l = getLeft();
        Date r = getRight();
        return String.format("[%s .. %s]", l.toString(), r.toString());
    }
}
\end{lstlisting}

\begin{itemize}
\tightlist
\item
  Fissa sia \passthrough{\lstinline!F!} che \passthrough{\lstinline!S!}
  al tipo concreto \passthrough{\lstinline!Date!}
  (\passthrough{\lstinline!extends Pair<Date, Date>!}).
\item
  Offre metodi con semantica di dominio più chiara:
  \passthrough{\lstinline!getLeft()!} e
  \passthrough{\lstinline!getRight()!}.
\end{itemize}

\paragraph{\texorpdfstring{\href{file:///home/wally/Documenti/GitHub/Programmazione-II/25-26/Teoria-e-Esercitazioni/Codice-20260902/Lezione14/MavenDate/src/main/java/it/oop/core/DateInterval.java}{DateInterval.java}}{DateInterval.java}}\label{dateinterval.java}

\begin{lstlisting}[language=Java]
package it.oop.core;

public class DateInterval extends OrderedPair<Date> {
    public DateInterval(Date left, Date right) {
        super(left, right);
    }

    public Date getLeft() { return getFirst(); }
    public Date getRight() { return getSecond(); }

    @Override
    public String toString() {
        Date l = getLeft();
        Date r = getRight();
        return String.format("[%s .. %s]", l.toString(), r.toString());
    }
}
\end{lstlisting}

\begin{itemize}
\tightlist
\item
  Eredita da \passthrough{\lstinline!OrderedPair<Date>!}: è legale
  perché \passthrough{\lstinline!Date implements Comparable<Date>!}.
\item
  Se le date passate non sono cronologicamente ordinate
  (\passthrough{\lstinline!left > right!}), il costruttore di
  \passthrough{\lstinline!OrderedPair!} emette il messaggio d'avviso
  \passthrough{\lstinline'"Pair not ordered!"'}.
\end{itemize}

\begin{center}\rule{0.5\linewidth}{0.5pt}\end{center}

\subsubsection{\texorpdfstring{2.5
\href{file:///home/wally/Documenti/GitHub/Programmazione-II/25-26/Teoria-e-Esercitazioni/Codice-20260902/Lezione14/MavenDate/src/main/java/it/oop/ui/MainDate.java}{MainDate.java}}{2.5 MainDate.java}}\label{maindate.java}

Raccoglie in un unico punto di ingresso tutte le feature viste nelle
lezioni 13 e 14:

\begin{lstlisting}[language=Java]
package it.oop.ui;

import it.oop.core.*;
import it.oop.core.Date;

public class MainDate {

    public static void main(String[] args) {
        // 1. Static Nested Class (Builder Pattern da Lezione 13)
        Date.Builder dateBuilder = new Date.Builder(2025);
        Date d1 = dateBuilder.build(10,9);
        Date d2 = dateBuilder.build(10,-1);
        System.out.println(d1.toString());
        System.out.println(d2.toString());

        // 2. Anonymous Class che implementa l'interfaccia Time (Lezione 13/14)
        Time init = new Time() {
            @Override public int getHours() { return 0; }
            @Override public int getMinutes() { return 0; }
            @Override public int getSeconds() { return 0; }
            @Override public String toString() {
                return String.format("%02d:%02d:%02d", getHours(), getMinutes(), getSeconds());
            }
        };
        System.out.println(init.toString());

        // 3. Lambda Expression che implementa FormattedDateConverter (Lezione 13/14)
        FormattedDateConverter toAmerican =
                d -> new AmericanDate(d.getDay(), d.getMonth(), d.getYear());
        System.out.println(toAmerican.convert(new ItalianDate(11, 11, 2025)) instanceof AmericanDate);

        // 4. Generics: Pair con Diamond Operator <>
        Pair<String, FormattedDate> event =
                new Pair<>("OOP exam", new ItalianDate(2, 2, 2026));
        System.out.println(event.getFirst() + " on " + event.getSecond().prettyPrint());

        // 5. DatePair e DateInterval
        DatePair dp = new DatePair(new Date(1,2,2025), new Date(1,1,2025));
        System.out.println("date pair: " + dp.toString());

        // Qui scatta l'avviso di ordinamento perché 1/2/2025 è successivo a 1/1/2025!
        DateInterval di = new DateInterval(new Date(1,2,2025), new Date(1,1,2025));

        // 6. Wildcard con Bounded Type (anticipazione Lezione 15)
        Pair<?, ? extends FormattedDate> unknown = event;
        System.out.println(unknown.getFirst().toString() + " on " + unknown.getSecond().prettyPrint());
    }
}
\end{lstlisting}

\begin{center}\rule{0.5\linewidth}{0.5pt}\end{center}

\subsection{3. Risultato di Compilazione ed Esecuzione
Reale}\label{risultato-di-compilazione-ed-esecuzione-reale}

Comando eseguito nel workspace del corso:

\begin{lstlisting}[language=bash]
mvn compile exec:java -Dexec.mainClass="it.oop.ui.MainDate"
\end{lstlisting}

Output ottenuto:

\begin{lstlisting}
y2025m9d10
y2025m1d1
00:00:00
true
OOP exam on 2 febbraio 2026
date pair: [y2025m2d1 .. y2025m1d1]
Pair not ordered!
OOP exam on 2 febbraio 2026
\end{lstlisting}

\subsubsection{Dettaglio riga per riga
dell'output:}\label{dettaglio-riga-per-riga-delloutput}

\begin{enumerate}
\def\labelenumi{\arabic{enumi}.}
\tightlist
\item
  \passthrough{\lstinline!y2025m9d10!} \(\rightarrow\) generato da
  \passthrough{\lstinline!d1 = dateBuilder.build(10, 9)!} (10 settembre
  2025).
\item
  \passthrough{\lstinline!y2025m1d1!} \(\rightarrow\) generato da
  \passthrough{\lstinline!d2 = dateBuilder.build(10, -1)!} (mese errato
  fallback su 1/1/2025).
\item
  \passthrough{\lstinline!00:00:00!} \(\rightarrow\) generato
  dall'override di \passthrough{\lstinline!toString()!} della
  \textbf{classe anonima} \passthrough{\lstinline!init!} che implementa
  \passthrough{\lstinline!Time!}.
\item
  \passthrough{\lstinline!true!} \(\rightarrow\) la \textbf{lambda}
  \passthrough{\lstinline!toAmerican!} converte l'istanza di
  \passthrough{\lstinline!ItalianDate!} in
  \passthrough{\lstinline!AmericanDate!}.
\item
  \passthrough{\lstinline!OOP exam on 2 febbraio 2026!} \(\rightarrow\)
  \passthrough{\lstinline!Pair<String, FormattedDate>!} con
  formattazione italiana.
\item
  \passthrough{\lstinline!date pair: [y2025m2d1 .. y2025m1d1]!}
  \(\rightarrow\) \passthrough{\lstinline!DatePair!} stampa le due date
  senza controllo di sequenzialità.
\item
  \passthrough{\lstinline"Pair not ordered!"} \(\rightarrow\) emesso dal
  costruttore di \passthrough{\lstinline!OrderedPair!} dentro
  \passthrough{\lstinline!new DateInterval(1/2/2025, 1/1/2025)!} perché
  \passthrough{\lstinline!1/2/2025 > 1/1/2025!}.
\item
  \passthrough{\lstinline!OOP exam on 2 febbraio 2026!} \(\rightarrow\)
  invocazione polimorfica tramite wildcard generica
  \passthrough{\lstinline!Pair<?, ? extends FormattedDate>!}.
\end{enumerate}

\begin{center}\rule{0.5\linewidth}{0.5pt}\end{center}

Dimmi \textbf{``vai''} per procedere con il \textbf{Blocco 14} (Lezione
15 / Slide T15 --- \emph{Generic Types}, Type Erasure, Wildcards, PECS e
diff in \passthrough{\lstinline!Lezione15/MavenDate!}).


\end{document}
